\documentclass{leveldoc}
\DocUses{foundations,slingshot,first_coil,rf_kick,ramp_the_coil}
\DocFingerprint{787b1465010b465a}

\begin{document}

\section{What the level asks}\label{sec:asks}

The game's text:
\leveldescription

One particle of charge $q = \val{q}$ and mass $m = \val{m}$ is launched from
$(\val{xZero}, \val{yZero})$ along $-x$ with $T_0 = \val{Tzero}$ in a world with $c = \val{c}$, and
the level includes radiation reaction: the particle loses the energy it radiates, by the
Landau--Lifshitz force (\physref{3.1}). A circular coil of radius $a = \val{a}$ and strength
$\kappa = \val{kappa}$ (field units, \physref{2.2}) surrounds the board's middle,
$(\val{xC}, \val{yC})$. The particle must enter the detector, the square
$\val{xDetLo} \le x \le \val{xDetHi}$, $\val{yDetLo} \le y \le \val{yDetHi}$ about the coil's
axis, before $t = \val{tMax}$. The player may place one magnet, a dipole perpendicular to the
board, on the coil's axis (the only node of the placement region), with a strength from the
listed magnitudes and either sign. A flight ends when the particle enters the detector, touches
the magnet's sphere (radius $\val{Rm}$) or the coil's wire (within $\val{wire}$ of its circle),
leaves the world, or reaches $t_{\max}$ (\xsecref{foundations}{sec:board}).

The reference places the magnet $\mu = \val{mu}$, whose field adds to the coil's near the axis.
The particle circles the axis, radiating, and spirals in; the stronger field makes it radiate
fast enough to reach the square at $t = \val{tRef}$. The document derives the reaction force in
this field, shows from the canonical angular momentum that without radiation no magnet lets the
particle near the axis, marches the flight with exact energy and torque balances as checks,
estimates the time by an adiabatic inspiral, and compares every legal placement with the game.

\section{Setup}\label{sec:setup}

The level file, field by field (the appendix reproduces it):
\setuptable

By \xeqref{foundations}{eq:launch} the particle starts with $\gamma_0 = \val{gammaZero}$, momentum
$p_\mathrm{i} = \val{pZero}$ and speed $v_0 = \val{vZero}$, at $\rho_0 = \val{rZero}$ from the
axis and moving perpendicular to the radius. The time scale of radiation reaction is
$\tau_0 = 2q^2/(3mc^3) = \val{tauZero}$. At the launch the field is $B_z = \val{BZero}$ (with the
magnet), the cyclotron frequency $\omega_\mathrm{c} = qB_z/(\gamma_0 m) = \val{omegaC}$ and the
period $\val{period}$. A circle about the axis through the launch would need the momentum
$q\rho_0B_z = \val{pCircle}$; the particle has a little less, so its orbit is a circle about a
point slightly off the axis, of mean radius $r_\mathrm{g} = \val{rGuide}$ (where
$q r_\mathrm{g} B_z(r_\mathrm{g}) = p_\mathrm{i}$). The detector's corners are
$\val{rCorner}$ from the axis, its edges $\val{rEdge}$.

\begin{figure}[htbp]
\centering
\begin{tikzpicture}
\begin{axis}[axis equal image, width=0.86\linewidth, xmin=-1.6, xmax=31.6, ymin=-1.6, ymax=21.6,
  xtick={0,5,...,30}, ytick={0,5,...,20}, tick label style={font=\small},
  xlabel={$x$}, ylabel={$y$}, label style={font=\small}]
  \draw[black!8] (0,0) grid (30,20);
  \draw[black!45, dashed] (-1,-1) rectangle (31,21);
  \fill[green!50!black, fill opacity=0.18] (\rawval{xDetLo},\rawval{yDetLo}) rectangle (\rawval{xDetHi},\rawval{yDetHi});
  \draw[green!40!black] (\rawval{xDetLo},\rawval{yDetLo}) rectangle (\rawval{xDetHi},\rawval{yDetHi});
  \draw[black!70, thick] (\rawval{xC},\rawval{yC}) circle[radius=\rawval{a}];
  \addplot[very thin, blue!70!black] table {pathref.dat};
  \fill[red!75!black] (\rawval{xC},\rawval{yC}) circle[radius=\rawval{Rm}];
  \fill (\rawval{xZero},\rawval{yZero}) circle[radius=0.15];
\end{axis}
\end{tikzpicture}
\caption{The coil (black circle), the magnet on its axis (red), the detector (green) and the
reference flight (blue), spiralling in from the launch (black dot) over about twelve turns.}
\label{fig:board}
\end{figure}

\section{Principles and model}\label{sec:model}

\begin{principles}
  \item The relativistic equation of motion \xeqref{foundations}{eq:eom} with the launch
    \xeqref{foundations}{eq:launch}, and the Landau--Lifshitz radiation reaction force
    (\LandauII{§§75--76}; \physref{3.1}, in the units of \xsecref{foundations}{sec:units}).
  \item Liénard's radiated power \xeqref{foundations}{eq:larmor}.
  \item The coil's field in its plane \xeqref{first_coil}{eq:loop-field} with
    \xeqref{first_coil}{eq:agm}, its vector potential \xeqref{ramp_the_coil}{eq:vector-potential};
    a magnetic dipole's field $\vect{B} = \mu\,(3(\uvec{m}\cdot\uvec{r})\uvec{r} - \uvec{m})/r^3$
    and vector potential $\vect{A} = \mu\,\uvec{m}\times\vect{r}/r^3$ (\Jackson{§5.6};
    \physref{2.2}); the derivatives of $K$ and $E$ (DLMF 19.4.1).
  \item The canonical angular momentum about a symmetry axis, as in
    \xeqref{ramp_the_coil}{eq:canonical} (\Jackson{§12.1}).
  \item The board's rules; the classical Runge--Kutta method, cubic Hermite interpolation and the
    error estimate from the step $2h$; Simpson's rule; bisection
    (\xsecref{foundations}{sec:board}, \xsecref{foundations}{sec:numerics}), as in
    \xsecref{slingshot}{sec:method}.
\end{principles}

The coil carries a steady current and the magnet is fixed: the field is static, purely magnetic
and symmetric about the axis. The particle starts at $\vect{x}(0) = (\val{xZero}, \val{yZero})$
with $\vect{p}(0) = (-p_\mathrm{i}, 0)$ and moves by \xeqref{foundations}{eq:eom} with the force
$q\vect{v}\times\vect{B} + \vect{f}$, $\vect{f}$ the reaction, until the first of the board's
events.

\begin{assumption}{The Landau--Lifshitz force}\label{as:ll}
The exact classical equation of a radiating point charge, the Lorentz--Abraham--Dirac equation,
has runaway solutions; the Landau--Lifshitz force is its reduction of order, exact to first
order in $\tau_0$ and free of runaways. It holds while the reaction is small against the Lorentz
force; here their ratio is $(2q^3B\gamma^2)/(3m^2c^3) = \val{ratio}$ at the launch, and that ratio
bounds the relative error of the reaction's effects (\physref{3.1}, test R6). The game uses the
same force, so this approximation does not enter the comparison.
\end{assumption}

\begin{assumption}{Fixed sources, a thin wire, a point particle}\label{as:fixed}
The coil's current and the magnet do not respond to the particle; the coil is a circle of zero
thickness whose wire only ends a flight (as in \levelref{first_coil}, its assumption
\xref{first_coil}{as:wire}); the particle is a point.
These are the game's model, and nothing here comes near the wire.
\end{assumption}

\section{Derivation}\label{sec:derivation}

\subsection{The field}\label{sec:field}

In the plane, at distance $\rho$ from the axis, the coil gives $b(\rho)$ of
\xeqref{first_coil}{eq:loop-field} and the dipole $\mu\uvec{z}$ gives
$-\mu/\rho^3$ ($\uvec{m}\cdot\uvec{r} = 0$ there), so
\begin{equation}\label{eq:field}
  B_z(\rho) = b(\rho) - \frac{\mu}{\rho^3},
  \qquad
  A_\phi(\rho) = a_\mathrm{coil}(\rho) + \frac{\mu}{\rho^2},
\end{equation}
with $a_\mathrm{coil}$ from \xeqref{ramp_the_coil}{eq:vector-potential}. A negative $\mu$ adds to
the coil's field, a positive one subtracts from it near the axis. The reaction needs the field's
gradient, radial by symmetry. Differentiating \xeqref{first_coil}{eq:loop-field},
$b = 2\kappa N/D$ with $N = (a^2 - \rho^2)E + \alpha^2K$ and $D = \alpha^2\beta$, by the chain rule
with $\dd m/\dd\rho = 4a(a - \rho)/\beta^3$,
\begin{equation}\label{eq:gradient}
  \deriv{K}{m} = \frac{E - (1 - m)K}{2m(1 - m)}, \qquad \deriv{E}{m} = \frac{E - K}{2m}, \qquad
  \deriv{b}{\rho} = 2\kappa\,\frac{N'D - ND'}{D^2},
\end{equation}
$N' = -2\rho E + (a^2 - \rho^2)E' + 2(\rho - a)K + \alpha^2K'$, $D' = 2(\rho - a)\beta + \alpha^2$;
the dipole adds $3\mu/\rho^4$. Central differences of $B_z$ confirm \eqref{eq:gradient} to
$\val{gradCheck}$ relative. With the reference magnet:
\begin{center}\small\val{fieldtab}\end{center}

\subsection{The reaction force in a static magnetic field}\label{sec:ll}

For the force $q(\vect{E} + \vect{v}\times\vect{B})$ the Landau--Lifshitz force is (\physref{3.1})
\begin{align*}
  \vect{f} &= \frac{2q^3}{3mc^3}\,\gamma\left[\frac{D\vect{E}}{Dt} + \vect{v}\times\frac{D\vect{B}}{Dt}\right] \\
  &\quad + \frac{2q^4}{3m^2c^4}\left[c\,\vect{E}\times\vect{B} + c\,\vect{B}\times(\vect{B}\times\vect{v})
    + \frac{\vect{E}(\vect{v}\cdot\vect{E})}{c}\right] \\
  &\quad - \frac{2q^4}{3m^2c^5}\,\gamma^2\vect{v}\left[(\vect{E} + \vect{v}\times\vect{B})^2
    - \frac{(\vect{E}\cdot\vect{v})^2}{c^2}\right],
\end{align*}
$D/Dt = \partial_t + \vect{v}\cdot\nabla$ along the world line. Here $\vect{E} = 0$, the field is
static and $\vect{B} = B_z\uvec{z}$ with $\vect{v}$ in the plane, so
$D\vect{B}/Dt = (\vect{v}\cdot\nabla B_z)\,\uvec{z}$,
$\vect{B}\times(\vect{B}\times\vect{v}) = -B_z^2\vect{v}$ and $(\vect{v}\times\vect{B})^2 = v^2B_z^2$;
with $1 + \gamma^2v^2/c^2 = \gamma^2$,
\begin{equation}\label{eq:ll-static}
  \vect{f} = \frac{2q^3}{3mc^3}\,\gamma\,(\vect{v}\cdot\nabla B_z)\,(\vect{v}\times\uvec{z})
  - \frac{2q^4B_z^2}{3m^2c^3}\,\gamma^2\,\vect{v}.
\end{equation}
The first term is perpendicular to $\vect{v}$; the second is a drag along it, growing with the
square of the field.

\subsection{Energy and angular momentum}\label{sec:balances}

Only the drag does work, so the energy falls at the rate
\begin{equation}\label{eq:balance}
  \deriv{}{t}\left(\gamma mc^2\right) = \vect{v}\cdot\vect{f} = -P,
  \qquad P = \frac{2q^4B_z^2}{3m^2c^3}\,\gamma^2v^2,
\end{equation}
and $T(t) + W(t) = T_0$ exactly, $W = \int_0^t P\,\dd t'$ the radiated energy. The Lorentz force
alone gives the acceleration $\vect{a}_\mathrm{L} = q\vect{v}\times\vect{B}/(\gamma m)$,
perpendicular to $\vect{v}$, and Liénard's power \xeqref{foundations}{eq:larmor} for it is
$\frac{2q^2}{3c^3}\gamma^6 a_\mathrm{L}^2(1 - v^2/c^2) = \frac{2q^2}{3c^3}\gamma^4\frac{q^2v^2B_z^2}{\gamma^2m^2}$:
the same $P$. With the full acceleration (the reaction's own included) Liénard's power differs at
the next order in $\tau_0$, the terms the Landau--Lifshitz reduction drops.

Without the reaction, the canonical angular momentum about the axis,
\begin{equation}\label{eq:canonical}
  J = (x - x_\mathrm{c})\,p_y - (y - y_\mathrm{c})\,p_x + q\,\rho A_\phi(\rho),
\end{equation}
is conserved: with $\vect{r}$ from the axis, $(\vect{r}\times q\vect{v}\times\vect{B})_z = -qB_z\,\rho\dot\rho$,
while $\dd(q\rho A_\phi)/\dd t = q\dot\rho\,\dd(\rho A_\phi)/\dd\rho = q\rho\dot\rho B_z$, as
$B_z = \rho^{-1}\dd(\rho A_\phi)/\dd\rho$. With it, only the reaction's torque changes $J$:
\begin{equation}\label{eq:torque}
  J(t) - J(0) = \int_0^t (\vect{r}\times\vect{f})_z\,\dd t'.
\end{equation}
\eqref{eq:balance} and \eqref{eq:torque} hold exactly for the model; the march checks both.

\subsection{Without radiation the axis is out of reach}\label{sec:bound}

Without the reaction, $\abs{\vect{p}} = p_\mathrm{i}$ (the magnetic force does no work) and $J$ is
conserved. At distance $\rho$ from the axis $\abs{(\vect{r}\times\vect{p})_z} \le \rho\,p_\mathrm{i}$,
so the particle can be there only if
\begin{equation}\label{eq:bound}
  \abs{J(0) - q\rho A_\phi(\rho)} \le \rho\,p_\mathrm{i}.
\end{equation}
Stepping $\rho$ inward from the launch and bisecting gives the smallest reachable distance, for
no magnet and every allowed one ($J$ at the launch in the second column):
\begin{center}\small\val{bounds}\end{center}
All of them stay farther than $\val{rBoundWorst}$ from the axis, outside the detector, whose
nearest points are $\val{rEdge}$ away: without radiation no placement solves the level. For the
reference the bound is $\val{rBoundRef}$, and a march without the reaction (at $h = 0.05$) ends
in a \val{noRadOut}, its least distance from the axis $\val{noRadRhoMin}$.

\subsection{The march}\label{sec:march}

There is no closed form. \xeqref{foundations}{eq:eom} with $\vect{f}$ of \eqref{eq:ll-static} is
marched by \xeqref{foundations}{eq:rk4} at $h = \val{H}$, with four more states: the radiated
energy $W$, Liénard's energy with the full and with the Lorentz acceleration, and the torque's
integral \eqref{eq:torque}. The events are located on the cubic Hermite interpolant by bisection
and the error is estimated against the step $2h$.

\subsection{The adiabatic inspiral}\label{sec:adiabatic}

The particle loses $\val{lossTurn}$\,\% of its kinetic energy in its first turn, and less in
later ones (the drag falls with the momentum); its orbit stays close to a circle about the axis
whose radius shrinks over about a dozen turns, and this model, which follows only the mean
radius, is judged by its result. On a
circle of radius $r$ about the axis the momentum is $p(r) = qrB_z(r)$ (its radius of curvature
$p/(qB_z)$ equals $r$), the energy $\mathcal{E}(r) = \sqrt{m^2c^4 + c^2p^2}$, and it falls at the
rate $P = 2q^4B_z^2p^2/(3m^4c^3)$ of \eqref{eq:balance} ($\gamma v = p/m$). So
$\dot r = -P/(\dd\mathcal{E}/\dd r)$ and the time from radius $r_1$ to $r_2$ is
\begin{equation}\label{eq:adiabatic}
  t = \int_{r_2}^{r_1} \frac{\dd\mathcal{E}/\dd r}{P(r)}\,\dd r,
  \qquad
  \deriv{\mathcal{E}}{r} = \frac{c^2p\,q\,(B_z + rB_z')}{\mathcal{E}}.
\end{equation}
From the mean radius $r_\mathrm{g} = \val{rGuide}$ to the detector's corners, $\val{rCorner}$,
Simpson's rule gives $t = \val{tAdCorner}$; to its edges, $\val{rEdge}$, $t = \val{tAdEdge}$. The
particle can enter between the two, depending on where on its orbit it is when it gets there.

\section{The reference solution by hand}\label{sec:hand}

\begin{procedure}
  \item\label{st:field} \emph{The field.} At each point, $\rho$ from the axis; $B_z$, its gradient
    and $A_\phi$ by \eqref{eq:field} and \eqref{eq:gradient}, as in the table of
    \S\ref{sec:field}; the force $q\vect{v}\times\vect{B} + \vect{f}$ with \eqref{eq:ll-static}.
  \item\label{st:march} \emph{The march.} \S\ref{sec:march} from the launch; every 100 time
    units (the last column the kinetic energy):

    \smallskip
    \centerline{\val{march}}
    \smallskip

    The particle enters the detector through its \val{faceRef} face at
    $(\val{xRef}, \val{yRef})$ at $t = \val{tRef}$, with $\vect{p} = (\val{pxRef}, \val{pyRef})$
    (step estimate $\val{errRef}$).
  \item\label{st:energy} \emph{The energy.} It radiated $W = \val{wLL}$, a fraction
    $\val{radFrac}$ of $T_0$, and arrives with $T = \val{tEnd}$; $T + W = T_0$ holds to
    $\val{wDev}$ along the march \eqref{eq:balance}. Liénard's energy with the Lorentz
    acceleration is $\val{wLienardL}$, the same; with the full acceleration $\val{wLienard}$,
    $\val{lorentzOnly}$ higher, the next order in $\tau_0$.
  \item\label{st:j} \emph{The angular momentum.} $J$ goes from $\val{jZero}$ to $\val{jEnd}$; the
    torque's integral is $\val{torque}$, and \eqref{eq:torque} holds to $\val{jDev}$.
  \item\label{st:adiabatic} \emph{The estimate.} The adiabatic inspiral \eqref{eq:adiabatic}
    reaches the detector's corners at $\val{tAdCorner}$; the march arrives $\val{adRel}$\,\% from
    that, entering near a corner.
\end{procedure}

\begin{result}[The reference flight]
With $\mu = \val{mu}$ the particle spirals in and enters the detector at
$(\val{xRef}, \val{yRef})$ at $t = \val{tRef} < t_{\max}$, having radiated $\val{radFrac}$ of its
kinetic energy. The level is solved.
\end{result}

\section{Comparison with the game}\label{sec:comparison}

\comparisontable

The flights last up to $\val{tMax}$, about fourteen turns, and the trajectories agree to about
$10^{-9}$ at their ends, the larger differences on the longest flights. The tolerances are ten
times the larger of the two computations' estimates or the game's preview--verification gap,
whichever is larger; for these long flights it is the gap (for the reference $\val{tolRef}$,
against ten times our estimate, $\val{errRefTen}$). The game reports Liénard's energy with the
acceleration of the Lorentz force alone, which \S\ref{sec:balances} shows equals the radiated
energy $W$ in this field; it agrees with ours to about $10^{-11}$ relative.

\section{Other solutions}\label{sec:others}

The table covers every legal placement: no magnet and each of the eight magnets. A magnet with
$\mu < 0$ strengthens the field near the axis; the drag \eqref{eq:ll-static} grows as $B_z^2$, and
the particle spirals in faster:
$\mu = -1$ arrives at $t = \val{tM10}$, $-2$ at $\val{tM20}$, $-4$ at $\val{tM40}$, having
radiated $\val{radM10}$, $\val{radM20}$ and $\val{radM40}$ of $T_0$. Without a magnet, or with
$\mu > 0$, which weakens the field near the axis, the particle radiates as much energy
($\val{radNone}$ of $T_0$ without a magnet, $\val{radP40}$ with $\mu = 4$) but its orbit shrinks
less, and at $t_{\max}$ it is still $\val{rhoEndNone}$ ($\mu = 0$), $\val{rhoEndP05}$, $\val{rhoEndP10}$,
$\val{rhoEndP20}$ and $\val{rhoEndP40}$ ($\mu = 0.5$ to $4$) from the axis: a \val{outNone} each.
On a circle $p = qrB_z$, so a weaker field near the axis means a smaller momentum, and the energy
must fall further before the orbit fits there. The reference is the weakest magnet that solves.

\section{What the game computes}\label{sec:engine}

The game evaluates the coil's and the magnet's fields in closed form (\physref{2.2}) and the
Landau--Lifshitz force in its general form (\physref{3.1}), with $D\vect{B}/Dt$ a central
difference along the world line over $10^{-5}$ cells (relative error about $10^{-10}$, far below
$\tau_0$'s order); in this field that is \eqref{eq:ll-static}. It integrates the motion with its
adaptive eighth-order Runge--Kutta method, the reaction's work as one more component
(\physref{5}), with event location and verification (\physref{6}, \physref{7}), and integrates
Liénard's power with the Lorentz force's acceleration by three-point Gauss--Legendre on each step.
The problem is the same as here; the differences in the comparison are the two integrations'
errors.

\section{Further reading}\label{sec:reading}

Radiation reaction and the Landau--Lifshitz force: \LandauII{§§75--76}, \Jackson{Ch.~16}.
Liénard's formula and synchrotron radiation: \Jackson{Ch.~14}. The coil and its vector
potential: \levelref{first_coil}, \levelref{ramp_the_coil}.

\levelfileappendix

\end{document}
